%! Unit 2, Lessons 2.1–2.5 — SAMPLE Test (blank)
% The practice form of the Unit 2 test. Item for item it is the parallel form of ../actual_test/:
% same two sections, the same ten multiple-choice targets in the same order, the same two
% free-response sets with the same parts and the same points (MC 10, FRQ 1 = 4, FRQ 2 = 6,
% total 20) — every context different, so working this one teaches the moves, not the answers.
% Merged into the unit student packet via ../../sample_test/ (see shared/tests.mk).
% NAME ROW: \namedateperiod IS carried — a unit test is taken in a testing setting, and the
% practice form is sat the same way (LESSON_SHAPE.md §4, components.md "Unit tests").
% NO SKETCH ITEMS: every display is pre-drawn in TikZ and read.
%
% THE DATA (every number checked; regression outputs computed from the listed points)
%  MC1  theater, 250: Under 30 84/16 (100), 30+ 90/60 (150); box office 76 -> 16/76 = 0.211.
%  MC2  bike share: commuters 240 (75% helmet = 180), recreational 60 (40% = 24).
%  MC4  14 stations, elevation (1000 ft) vs July mean temp: r = -0.95, one station +10.3 off.
%  MC6  yhat = 4.0 + 2.6(200) = 524 cm (data 10–60 days).
%  MC7  yhat = 1.5 + 0.9(12) = 12.3; residual 11.2 - 12.3 = -1.1 lb.
%  MC9  22 runners, miles x vs 5K minutes: a 31.64, b -0.2103, S 1.396, R-Sq 76.1% (r -0.872).
%  MC10 ramp, t 0.5..3.0 s: yhat = -48.5 + 69.9t, r = 0.98, residuals
%       19.0 7.0 -0.8 -8.1 -11.1 -11.8 -11.7 -7.3 -1.9 8.1 18.6 (a U).
%  FRQ1 college, 400: FT 132/84/24 (240) = 55/35/10%; PT 112/36/12 (160) = 70/22.5/7.5%.
%  FRQ2 15 refrigerators, capacity x (cu ft) vs energy y (kWh/yr):
%    x 10.1 11.5 12.8 14.0 15.2 16.4 17.9 18.6 19.8 20.5 21.9 23.0 24.7 25.6 27.3
%    y  379  429  391  459  444  461  539  504  516  546  577  563  603  576  616
%    a 249.70, b 13.811, S 21.76, R-Sq 92.4% (r 0.961); (17.9, 539): yhat 496.9, res +42.1.
%
% PAGE LOCKSTEP with ../../test_keys/practice_test_key/: every \writelines{n} here is answered
% by exactly n \ansline{} there, each terse enough not to wrap; every work block is
% byte-identical. Nothing else differs except the preamble, the header, and the \ans{} marks
% on the correct MC options.
\documentclass[12pt]{article}
\usepackage{apstats-article}
\usepackage{apstats-boxes}

\begin{document}

\pageheader{Unit 2 \quad Lessons 2.1--2.5}{Sample Test}

\namedateperiod

{\small\hfill Score:\ \blank{1.4cm}\ / 20}

\vspace{0.06in}

\boxguard[8]
\begin{remindbox}
\small
This is the \textbf{practice} test --- not scored. Work it in one period, notes closed, then check
it against the key. It matches the real test item for item; every answer must be in context.
\end{remindbox}

\vspace{0.10in}

\boxguard[16]
\begin{headlinebox}{goldbg}
\small
\textbf{Section I --- Multiple Choice} \hfill \textbf{10 points, 1 each}\\
Circle the single best answer. Read all five options --- more than one of them is about the
context, not just the arithmetic.
\end{headlinebox}

\vspace{0.10in}

\begin{enumerate}[leftmargin=1.9em, itemsep=0.13in]

  \boxguard[12]
  \item \begin{minipage}[t]{0.44\linewidth}\raggedright
        A community theater asked 250 audience members their age and how they bought their
        ticket. Of the audience members who bought at the \textbf{box office}, what proportion
        were under 30?
        \end{minipage}\hfill
        \begin{minipage}[t]{0.52\linewidth}
        \centering\small
        \renewcommand{\arraystretch}{1.15}
        \begin{tabular}[t]{@{} l | c c | c @{}}
         & \textbf{Online} & \textbf{Box office} & \textbf{Total} \\ \hline
        \textbf{Under 30} & 84 & 16 & 100 \\
        \textbf{30 and over} & 90 & 60 & 150 \\ \hline
        \textbf{Total} & 174 & 76 & 250 \\
        \end{tabular}
        \end{minipage}
        \begin{enumerate}[label=\textbf{(\Alph*)}, leftmargin=2.2em, itemsep=2pt]
          \item $16/250 = 0.064$
          \item $16/100 = 0.160$
          \item $16/76 \approx 0.211$
          \item $76/250 = 0.304$
          \item $100/250 = 0.400$
        \end{enumerate}

  \boxguard[15]
  \item \begin{minipage}[t]{0.62\linewidth}\raggedright
        A city bike-share program surveyed 300 riders: 240 who ride to commute and 60 who ride
        for recreation. The segmented bar graph shows whether each group wore a helmet on
        their last ride. Which statement is best supported?
        \end{minipage}\hfill
        \begin{minipage}[t]{0.34\linewidth}
        \centering
        \begin{tikzpicture}[x=1cm, y=0.024cm, baseline=(current bounding box.north)]
          \foreach \y in {0,50,100} {
            \draw (-0.08,\y) -- (0.08,\y) node[left, font=\tiny, xshift=-2pt] {\y\%};}
          \fill[navy]    (0.3,0)  rectangle (1.2,75);
          \fill[goldacc] (0.3,75) rectangle (1.2,100);
          \fill[navy]    (1.6,0)  rectangle (2.5,40);
          \fill[goldacc] (1.6,40) rectangle (2.5,100);
          \draw (0.3,0) rectangle (1.2,100);
          \draw (1.6,0) rectangle (2.5,100);
          \draw[thick] (0,0) -- (0,104);
          \draw[thick] (0,0) -- (2.7,0);
          \node[below, font=\tiny, align=center] at (0.75,-2) {Commute\\(240)};
          \node[below, font=\tiny, align=center] at (2.05,-2) {Recreation\\(60)};
          \fill[navy] (0.2,112) rectangle (0.4,122); \node[right, font=\tiny] at (0.4,117) {Helmet};
          \fill[goldacc] (1.4,112) rectangle (1.6,122); \node[right, font=\tiny] at (1.6,117) {No helmet};
        \end{tikzpicture}
        \end{minipage}
        \begin{enumerate}[label=\textbf{(\Alph*)}, leftmargin=2.2em, itemsep=2pt]
          \item No association: far more commuters than recreational riders were surveyed.
          \item Commuters skip the helmet more often: 60 commuters rode without one, but only
                36 recreational riders did.
          \item There is no association, because both bars reach 100\%.
          \item Type of ride and helmet use are associated: 75\% of commuters wore a helmet,
                compared with 40\% of recreational riders.
          \item Commuting causes riders to wear a helmet.
        \end{enumerate}

  \boxguard[9]
  \item A nutritionist wants to know whether the \textbf{sugar} in a breakfast cereal (grams per
        serving) helps predict its \textbf{customer rating} (0 to 100). Which is correct?
        \begin{enumerate}[label=\textbf{(\Alph*)}, leftmargin=2.2em, itemsep=2pt]
          \item Explanatory: sugar per serving. \ Response: customer rating.
          \item Explanatory: customer rating. \ Response: sugar per serving.
          \item Both are response variables, because the nutritionist assigned neither one.
          \item The cereals are the explanatory variable; the ratings are the response.
          \item There is no explanatory variable, because this is not an experiment.
        \end{enumerate}

  \boxguard[13]
  \item \begin{minipage}[t]{0.50\linewidth}\raggedright
        A weather service recorded the \textbf{elevation} (thousands of feet) and the average
        \textbf{July high temperature} ($^\circ$F) at 14 weather stations. Which description of
        the scatterplot is best?
        \end{minipage}\hfill
        \begin{minipage}[t]{0.46\linewidth}
        \centering
        \begin{tikzpicture}[x=0.55cm, y=0.085cm, baseline=(current bounding box.north)]
          \foreach \y in {50,60,...,90} {
            \draw[linegray, very thin] (0,\y) -- (10.5,\y);
            \node[left, font=\tiny] at (0,\y) {\y};}
          \draw[thick] (0,48) -- (10.7,48);
          \draw[thick] (0,48) -- (0,91);
          \foreach \x in {0,2,...,10} {
            \draw (\x,47.3) -- (\x,48.7); \node[below, font=\tiny] at (\x,47.5) {\x};}
          \foreach \x/\y in {1/84.9, 2/81.5, 2.5/79.0, 3/77.8, 4/74.2, 4.5/72.9, 5/70.1,
                             5.5/80.0, 6/67.4, 6.5/64.9, 7/63.8, 8/60.1, 9/56.2, 10/53.5}
            \fill[navy] (\x,\y) circle (1.7pt);
          \node[font=\tiny] at (5.2,40) {Elevation (1000 ft)};
          \node[rotate=90, font=\tiny] at (-1.5,69.5) {July high ($^\circ$F)};
        \end{tikzpicture}
        \end{minipage}
        \begin{enumerate}[label=\textbf{(\Alph*)}, leftmargin=2.2em, itemsep=2pt]
          \item Strong, positive, and linear, with no unusual points.
          \item Strong, negative, and linear, with one station far above the pattern.
          \item Weak, negative, and curved, with no unusual points.
          \item Moderate, positive, and curved, with one station far below the pattern.
          \item No association, because one station does not follow the pattern.
        \end{enumerate}

  \boxguard[10]
  \item For 40 commuters, the correlation between one-way commute \textbf{distance} (miles) and
        monthly \textbf{fuel cost} (dollars) is $r = 0.64$. Which statement is true?
        \begin{enumerate}[label=\textbf{(\Alph*)}, leftmargin=2.2em, itemsep=2pt]
          \item If fuel cost were recorded in euros instead of dollars, $r$ would change.
          \item If fuel cost were put on the $x$-axis and distance on the $y$-axis, $r$ would be
                $-0.64$.
          \item $r = 0.64$ means 64\% of the variation in fuel cost is explained by distance.
          \item Had $r$ been close to 0, distance and fuel cost could have no relationship of
                any kind.
          \item If distance were recorded in kilometers instead of miles, $r$ would still be
                0.64.
        \end{enumerate}

  \boxguard[8]
  \item A gardener's line predicts a sunflower's height (cm) from its age (days):
        $\widehat{\text{height}} = 4.0 + 2.6(\text{age})$, built from plants 10 to 60 days old.
        What does the line predict for a sunflower \textbf{200 days} old, and should the
        prediction be trusted?
        \begin{enumerate}[label=\textbf{(\Alph*)}, leftmargin=2.2em, itemsep=2pt]
          \item 524 cm, and it can be trusted, because the least-squares line is the best line.
          \item 520 cm, and it can be trusted, because the intercept matters only for young
                plants.
          \item 524 cm, but 200 days is far outside the 10--60 days in the data, so this
                extrapolation may be badly wrong.
          \item 802.6 cm, but 200 days is an extrapolation.
          \item No prediction is possible, because no 200-day-old sunflower was measured.
        \end{enumerate}

  \boxguard[9]
  \item A veterinarian's line predicts a puppy's weight (lb) from its age (weeks), built from
        puppies 6 to 20 weeks old: $\widehat{\text{weight}} = 1.5 + 0.9(\text{age})$. A
        12-week-old puppy weighs 11.2 lb. What is its residual, and what does it mean?
        \begin{enumerate}[label=\textbf{(\Alph*)}, leftmargin=2.2em, itemsep=2pt]
          \item $1.1$ lb; the puppy weighs 1.1 lb more than the line predicts.
          \item $-1.1$ lb; the puppy weighs 1.1 lb less than the line predicts.
          \item $-1.1$ lb; the puppy weighs 1.1 lb more than the line predicts.
          \item $1.1$ lb; the line predicted a weight 1.1 lb too low.
          \item $12.3$ lb; the residual is the predicted weight.
        \end{enumerate}

  \boxguard[11]
  \item For 32 schools, a district fit
        $\widehat{\text{cost}} = 14.2 + 0.85(\text{area})$, where cost is the annual electricity
        cost (thousands of dollars) and area is the floor area (thousands of square feet).
        Which is the correct interpretation of the slope?
        \begin{enumerate}[label=\textbf{(\Alph*)}, leftmargin=2.2em, itemsep=2pt]
          \item Every school's electricity cost rises by exactly \$850 for each extra
                1{,}000 square feet.
          \item A school with 0 square feet is predicted to spend \$14{,}200 a year.
          \item Adding 1{,}000 square feet to a school will cause its annual cost to rise by
                \$850.
          \item For each additional 1{,}000 square feet of floor area, the \emph{predicted}
                annual electricity cost increases by \$850.
          \item About 85\% of the variation in electricity cost is explained by floor area.
        \end{enumerate}

  \boxguard[12]
  \item \begin{minipage}[t]{0.42\linewidth}\raggedright
        A coach regressed 5K race time (minutes) on weekly training miles for 22 runners.
        Which is the correct interpretation of R-Sq?
        \end{minipage}\hfill
        \begin{minipage}[t]{0.56\linewidth}
        \centering\footnotesize\ttfamily
        \begin{tabular}[t]{@{} l r r r r @{}}
        Predictor & Coef & SE Coef & T & P \\ \hline
        Constant & 31.645 & 0.714 & 44.33 & 0.000 \\
        Miles & $-$0.2103 & 0.0264 & $-$7.97 & 0.000 \\
        \end{tabular}\par
        S = 1.396 \quad R-Sq = 76.1\% \quad R-Sq(adj) = 74.9\%\par
        \end{minipage}
        \begin{enumerate}[label=\textbf{(\Alph*)}, leftmargin=2.2em, itemsep=2pt]
          \item About 76.1\% of the variation in 5K time is explained by the linear model with
                weekly training miles.
          \item The correlation between training miles and 5K time is 0.761.
          \item The line predicts 76.1\% of the runners' 5K times correctly.
          \item Each additional training mile lowers a runner's 5K time by 76.1\%.
          \item Weekly training miles cause 76.1\% of a runner's 5K time.
        \end{enumerate}

  \boxguard[14]
  \item \begin{minipage}[t]{0.50\linewidth}\raggedright
        A physics class rolled a ball down a ramp and recorded the \textbf{distance} it had
        traveled (cm) at 11 \textbf{times} (seconds). The least-squares line is
        $\hat{y} = -48.5 + 69.9x$, with $r = 0.98$. Its residual plot is shown. Which
        conclusion is best?
        \end{minipage}\hfill
        \begin{minipage}[t]{0.46\linewidth}
        \centering
        \begin{tikzpicture}[x=1.75cm, y=0.068cm, baseline=(current bounding box.north)]
          \draw[linegray] (0.25,0) -- (3.2,0);
          \draw[thick] (0.25,-16) -- (0.25,22);
          \draw[thick] (0.25,-16) -- (3.2,-16);
          \foreach \y in {-10,0,10,20} { \node[left, font=\tiny] at (0.25,\y) {\y}; }
          \foreach \x in {0.5,1.0,...,3.0} {
            \draw (\x,-16.7) -- (\x,-15.3); \node[below, font=\tiny] at (\x,-16.3) {\x};}
          \foreach \x/\y in {0.5/19.0, 0.75/7.0, 1.0/-0.8, 1.25/-8.1, 1.5/-11.1, 1.75/-11.8,
                             2.0/-11.7, 2.25/-7.3, 2.5/-1.9, 2.75/8.1, 3.0/18.6}
            \fill[navy] (\x,\y) circle (1.7pt);
          \node[font=\tiny] at (1.72,-27) {Time (seconds)};
          \node[rotate=90, font=\tiny] at (-0.05,3) {Residual (cm)};
        \end{tikzpicture}
        \end{minipage}
        \begin{enumerate}[label=\textbf{(\Alph*)}, leftmargin=2.2em, itemsep=2pt]
          \item The line is a good model, because $r = 0.98$ is very close to 1.
          \item The line is a good model, because the residuals add to zero.
          \item Time and distance are not associated, because the residuals are not random.
          \item The points at 0.5 s and 3.0 s are influential and should be removed.
          \item A line is not an appropriate model: the residuals form a clear U-shape, so the
                form is curved even though $r$ is high.
        \end{enumerate}

\end{enumerate}

\vspace{0.14in}

% No \newpage here: the guard keeps the Section II strip off a page edge.
\boxguard[30]
\begin{headlinebox}{goldbg}
\small
\textbf{Section II --- Free Response} \hfill \textbf{10 points}\\
Show your reasoning. Answer \emph{in context} --- name the variables and their units.
\end{headlinebox}

\vspace{0.10in}

\begin{enumerate}[leftmargin=1.9em, itemsep=0.16in, resume]

  \boxguard[22]
  \item \textbf{(4 points)} A community college asked 400 students whether they attend
        \textbf{full-time} or \textbf{part-time} and how they \textbf{usually get to campus}.

        \vspace{4pt}
        \begin{minipage}[c]{0.56\linewidth}
        \centering\small
        \setlength{\tabcolsep}{4pt}\renewcommand{\arraystretch}{1.15}
        \begin{tabular}{@{} l | c c c | c @{}}
         & \textbf{Drive} & \textbf{Bus} & \textbf{Walk/bike} & \textbf{Total} \\ \hline
        \textbf{Full-time} & 132 & 84 & 24 & 240 \\
        \textbf{Part-time} & 112 & 36 & 12 & 160 \\ \hline
        \textbf{Total} & 244 & 120 & 36 & 400 \\
        \end{tabular}
        \end{minipage}\hfill
        \begin{minipage}[c]{0.40\linewidth}
        \centering
        \begin{tikzpicture}[x=1cm, y=0.026cm]
          \foreach \y in {0,50,100} {
            \draw (-0.08,\y) -- (0.08,\y) node[left, font=\tiny, xshift=-2pt] {\y\%};}
          \fill[navy]     (0.3,0)   rectangle (1.2,55);
          \fill[goldacc]  (0.3,55)  rectangle (1.2,90);
          \fill[skymid]   (0.3,90)  rectangle (1.2,100);
          \fill[navy]     (1.6,0)   rectangle (2.5,70);
          \fill[goldacc]  (1.6,70)  rectangle (2.5,92.5);
          \fill[skymid]   (1.6,92.5) rectangle (2.5,100);
          \draw (0.3,0) rectangle (1.2,100);
          \draw (1.6,0) rectangle (2.5,100);
          \node[white, font=\tiny\bfseries] at (0.75,27.5) {55\%};
          \node[font=\tiny\bfseries] at (0.75,72.5) {35\%};
          \node[white, font=\tiny\bfseries] at (2.05,35) {70\%};
          \node[font=\tiny\bfseries] at (2.05,81.25) {22.5\%};
          \draw[thick] (0,0) -- (0,104);
          \draw[thick] (0,0) -- (2.7,0);
          \node[below, font=\tiny] at (0.75,-2) {Full-time};
          \node[below, font=\tiny] at (2.05,-2) {Part-time};
          \fill[navy] (3.0,85) rectangle (3.2,95); \node[right, font=\tiny] at (3.2,90) {Drive};
          \fill[goldacc] (3.0,60) rectangle (3.2,70); \node[right, font=\tiny] at (3.2,65) {Bus};
          \fill[skymid] (3.0,35) rectangle (3.2,45); \node[right, font=\tiny] at (3.2,40) {Walk/bike};
        \end{tikzpicture}
        \end{minipage}

        \vspace{6pt}
        \textbf{(a)} What proportion of \emph{all} 400 students ride the bus? Of the students who
        walk or bike, what proportion attend full-time?\\[4pt]
        \writelines{2}

        \boxguard[6]
        \vspace{6pt}
        \textbf{(b)} Use the graph to compare how full-time and part-time students get to campus.\\[4pt]
        \writelines{2}

        \boxguard[6]
        \vspace{6pt}
        \textbf{(c)} Are enrollment status and travel to campus associated? Justify with numbers.
        Does this survey show that attending part-time \emph{causes} a student to drive?\\[4pt]
        \writelines{3}

  \boxguard[26]
  \item \textbf{(6 points)} An appliance lab measured the \textbf{capacity} (cubic feet) and the
        \textbf{annual energy use} (kWh per year) of 15 refrigerator models, then regressed energy
        use on capacity.

        \vspace{4pt}
        \noindent\hfil
        \begin{tikzpicture}[x=0.26cm, y=0.0135cm, baseline=(current bounding box.south)]
          \foreach \y in {400,500,600} {
            \draw[linegray, very thin] (8,\y) -- (29,\y);
            \node[left, font=\tiny] at (8,\y) {\y};}
          \draw[thick] (8,340) -- (29.3,340);
          \draw[thick] (8,340) -- (8,660);
          \foreach \x in {8,12,...,28} {
            \draw (\x,335) -- (\x,345); \node[below, font=\tiny] at (\x,337) {\x};}
          \foreach \x/\y in {10.1/379, 11.5/429, 12.8/391, 14.0/459, 15.2/444, 16.4/461,
                             17.9/539, 18.6/504, 19.8/516, 20.5/546, 21.9/577, 23.0/563,
                             24.7/603, 25.6/576, 27.3/616}
            \fill[navy] (\x,\y) circle (1.7pt);
          \node[font=\tiny] at (18.5,290) {Capacity (cubic feet)};
          \node[rotate=90, font=\tiny] at (3.2,500) {Energy use (kWh/yr)};
          \node[font=\bfseries\scriptsize] at (18.5,690) {Scatterplot};
        \end{tikzpicture}\hfil
        \begin{tikzpicture}[x=0.26cm, y=0.052cm, baseline=(current bounding box.south)]
          \draw[linegray] (8,0) -- (29,0);
          \draw[thick] (8,-45) -- (29.3,-45);
          \draw[thick] (8,-45) -- (8,48);
          \foreach \y in {-40,-20,0,20,40} { \node[left, font=\tiny] at (8,\y) {\y}; }
          \foreach \x in {8,12,...,28} {
            \draw (\x,-46.3) -- (\x,-43.7); \node[below, font=\tiny] at (\x,-45.5) {\x};}
          \foreach \x/\y in {10.1/-10.2, 11.5/20.5, 12.8/-35.5, 14.0/15.9, 15.2/-15.6,
                             16.4/-15.2, 17.9/42.1, 18.6/-2.6, 19.8/-7.2, 20.5/13.2,
                             21.9/24.8, 23.0/-4.4, 24.7/12.2, 25.6/-27.3, 27.3/-10.8}
            \fill[navy] (\x,\y) circle (1.7pt);
          \node[font=\tiny] at (18.5,-58) {Capacity (cubic feet)};
          \node[rotate=90, font=\tiny] at (3.2,0) {Residual (kWh/yr)};
          \node[font=\bfseries\scriptsize] at (18.5,56) {Residual plot};
        \end{tikzpicture}\hfil

        \vspace{4pt}
        {\centering\footnotesize\ttfamily
        \begin{tabular}{@{} l r r r r @{}}
        Predictor & Coef & SE Coef & T & P \\ \hline
        Constant & 249.70 & 21.23 & 11.76 & 0.000 \\
        Capacity & 13.811 & 1.100 & 12.56 & 0.000 \\
        \end{tabular}\par
        S = 21.76 \quad R-Sq = 92.4\% \quad R-Sq(adj) = 91.8\%\par}

        \boxguard[5]
        \vspace{6pt}
        \textbf{(a)} Describe the association between capacity and annual energy use.\\[4pt]
        \writelines{2}

        \boxguard[6]
        \vspace{6pt}
        \textbf{(b)} Write the equation of the least-squares regression line, using the variable
        names. Then interpret the slope in context.\\[4pt]
        \writelines{3}

        \boxguard[12]
        \vspace{6pt}
        \textbf{(c)} One model has a capacity of 17.9 cubic feet and uses 539 kWh per year. Find
        its residual and interpret it in context.
        \begin{work}
          \widehat{\text{energy}} &= 249.70 + 13.811(17.9) \\
                                  &\approx 496.9 \text{ kWh} \\
          \text{residual} &= y - \hat{y} \\
                          &= 539 - 496.9 \\
                          &= 42.1 \text{ kWh}
        \end{work}
        \writelines{2}

        \boxguard[5]
        \vspace{6pt}
        \textbf{(d)} Interpret the value of $r^2$ in context.\\[4pt]
        \writelines{2}

        \boxguard[5]
        \vspace{6pt}
        \textbf{(e)} Is a linear model appropriate for these data? Justify your answer using the
        residual plot.\\[4pt]
        \writelines{2}

\end{enumerate}

\end{document}
